% =============================================================
%  ch04_stereosehen.tex  --  Stereosehen und Merkmalserkennung
% =============================================================

\chapter{Stereosehen und Merkmalserkennung}
\label{ch:stereosehen}

\section{Lochkameramodell und Tiefe aus Disparität}
Eine kalibrierte Kamera bildet einen 3D-Punkt über die Intrinsik-Matrix $\bm{K}$
ab:
\begin{equation}
\bm{K} = \begin{pmatrix} f_x & 0 & c_x \\ 0 & f_y & c_y \\ 0 & 0 & 1 \end{pmatrix}
\end{equation}
Bei zwei horizontal um die \textbf{Basislinie} $b$ versetzten Kameras erscheint
ein Punkt im linken und rechten Bild an leicht verschiedener Spalte; die
Differenz ist die \textbf{Disparität} $d$. Die Tiefe folgt direkt:
\begin{equation}
\boxed{\,Z = \dfrac{f\cdot b}{d}\,}
\end{equation}
Große Basislinie $\rightarrow$ bessere Tiefenauflösung in der Ferne, aber
kleinerer überlappender Sichtbereich in der Nähe. Die Disparität wird auf
\emph{rektifizierten} Bildern berechnet (beide Bildebenen koplanar,
Epipolarlinien horizontal), typisch mit \code{StereoBM} (schnell) oder
\code{StereoSGBM} (genauer) aus OpenCV.

\section{Kalibrierung}
Beide Kameras müssen kalibriert werden (Intrinsik + Verzerrung) und zueinander
(Extrinsik: $\bm{R},\bm{t}$ zwischen den Linsen). Standard ist ein
Schachbrettmuster \cite{opencv_calib} und das ROS\,2-Paket
\code{camera\_calibration} aus dem
\code{image\_pipeline}. Ergebnis sind \code{.yaml}-Dateien, die der Treiber
lädt. \textbf{Ohne saubere Kalibrierung ist jede Tiefenmessung wertlos} -- das
ist der erste reale Stolperstein gegenüber der Simulation, wo die ``Kamera''
perfekt kalibriert ist. Schon kleine Kalibrierfehler skalieren systematisch in
die Objektgenauigkeit hinein \cite{kahmen2020stereocalib}; die praktische
Durchführung samt detektiertem Schachbrettmuster ist in
Kapitel~\ref{ch:kalibrierung} ausgeführt.

\section{Merkmalserkennung ohne eigenes KI-Training}
Es gibt ausgereifte \emph{klassische} Detektoren, die ohne Trainingsdaten und
ohne neuronale Netze auskommen. Sie finden markante, wiedererkennbare
Bildpunkte (Ecken, Blobs) und beschreiben deren lokale Umgebung:
\begin{longtable}{@{}p{2.4cm}p{4.5cm}p{8cm}@{}}
\toprule
\textbf{Detektor} & \textbf{Eigenschaft} & \textbf{Einsatz} \\
\midrule
\textbf{ORB}   & schnell, binär, lizenzfrei    & Echtzeit-Matching, VO, Standardwahl für Hobby \\
\textbf{SIFT}  & skalierungs-/rotationsinvariant, robust & genaues Matching, CAD-Posenschätzung \\
\textbf{AKAZE} & gute Balance                  & Alternative zu SIFT/ORB \\
\textbf{BRISK} & binär, schnell                & ressourcenarme Plattformen \\
\bottomrule
\end{longtable}
\noindent Alle sind in OpenCV enthalten (\code{cv2.ORB\_create()} usw.) und über
\code{cv\_bridge} in ROS\,2 nutzbar. Das Matching erfolgt mit \code{BFMatcher}
(Brute Force) oder \code{FLANN}, gefolgt von Ausreißerfilterung
(Lowe-Ratio-Test, RANSAC). Eine quantitative Gegenüberstellung von Genauigkeit
und Geschwindigkeit dieser Detektoren findet sich in
Kapitel~\ref{ch:merkmalserkennung} \cite{tareen2018comparative}. Für \textbf{Bewegungserkennung} (Projekt~3) brauchst
du nicht einmal Merkmale: Hintergrundsubtraktion
(\code{cv2.createBackgroundSubtractorMOG2}) oder optischer Fluss
(\code{calcOpticalFlowFarneback}) reichen.

\section{Posenschätzung -- 6D-Lage eines Bauteils}
Für Projekt~4 (CAD-Matching) gibt es zwei Hauptwege:
\begin{enumerate}
\item \textbf{2D-3D über PnP:} Merkmale im Kamerabild gegen bekannte 3D-Punkte
  des CAD-Modells matchen, dann \code{cv2.solvePnPRansac()} liefert Rotation +
  Translation des Bauteils relativ zur Kamera.
\item \textbf{3D-3D über ICP:} Aus der Stereokamera eine Punktwolke erzeugen und
  per \textbf{Iterative Closest Point} an die aus dem CAD-Modell gesampelte
  Punktwolke anpassen (z.\,B.\ mit Open3D). Liefert die volle 6D-Pose, robust
  bei guter Initialschätzung.
\end{enumerate}
Das Ergebnis (Rotation als Quaternion + Translation) wird in TF veröffentlicht
und vom Greifplaner (MoveIt\,2) konsumiert.
